\section{Incidence, degree sequences, and walk counts}\label{sec:basic}
An undirected multigraph on $[n]$ with $m$ labelled edges has two endpoint
tables $s_e,t_e<n$, $e<m$. Equality of the endpoints gives a loop. An edge
with identical endpoint data but a different label is still a different edge.
For a directed graph the order specifies tail and head. A simple graph has no
loops and at most one edge on each unordered pair. Its symmetric zero-diagonal
adjacency bit matrix is $A$. The upper triangle, in the order
$(0,1),(0,2),(1,2),(0,3),\ldots$, uses one bit per pair. This encoding and the
square encoding in the Lean file are interconvertible by bounded digit scans.
A mixed graph marks each record as directed or undirected. Natural weights
are a further digit table, with maximum found by a scan. A binary relation on
$[n]$ is exactly the arc mask of a digraph permitting a loop at each vertex.

\begin{proposition}[Incidence identities and parity]\leavevmode\label{prop:degree}
Put $d(v)=\sum_{e<m}([s_e=v]+[t_e=v])$, where brackets denote truth bits.
Then $\sum_vd(v)=2m$, and the number of odd-degree vertices is even. In a
digraph the indegrees and outdegrees each sum to $m$. For every cut $S$,
$\sum_{v\in S}d(v)=2|E(S)|+|\delta(S)|$, with $\delta(S)$ the edges with
exactly one end in $S$. In an $r$-regular graph $nr=2m$.
\end{proposition}
\begin{proof}
Fix an edge. Summing each endpoint indicator over all vertices gives one,
including when both endpoints coincide. Summing over a subset gives two,
one or zero according as both, one or neither endpoint lies there. Exchange
the finite sums. Write $d(v)=2q_v+b_v$ with $b_v\in\{0,1\}$. Reducing the
sum modulo two proves the parity assertion. Keeping only tail or head indicators
gives the directed identities. Substitution gives the regular case.
\end{proof}

\begin{proposition}[Simple graph degree bounds]\leavevmode\label{idx:elementary-simple-graph-degree-bounds}
For $n\ge 1$, $0\le d(v)\le n-1$ and $m\le\binom n2$. Equality in the
edge bound characterizes $K_n$. For $n\ge2$ some two vertices have the same
degree. Complementation replaces each degree $d$ by $n-1-d$.
\end{proposition}
\begin{proof}
Row $v$ has $n-1$ possible off-diagonal entries and the upper triangle has
$\sum_{j<n}j=\binom n2$ slots. A missing slot makes the edge sum strictly
smaller. A row of degree $n-1$ meets every other vertex, excluding an isolated
vertex. Scan for such a row. On success the degrees lie in $[1,n-1]$ and
otherwise in $[0,n-2]$. Either range has $n-1$ values. Apply the finite
pigeonhole argument. Complementation flips every off-diagonal slot.
\end{proof}

\begin{theorem}[Havel--Hakimi realization criterion]\leavevmode\label{idx:elementary-havel-hakimi-realization-criterion}
Let $d_0\ge\cdots\ge d_{n-1}\ge0$ be labelled demands for a simple graph.
The empty list is realizable. For a nonempty list set $d=d_0$. Reject if
$d\ge n$ or one of $d_1,\ldots,d_d$ is zero. Otherwise remove $d_0$,
subtract one from $d_1,\ldots,d_d$, and sort. The initial list is realizable
if and only if the reduced list is. The sequence of successful reductions
reconstructs a realizing graph.
\end{theorem}
\begin{proof}
In a realization write $v=0$ and $T=\{1,\ldots,d\}$. If some $x\in T$
misses $v$, there is a neighbour $y$ of $v$ outside $T$ because $v$ has
$d$ distinct neighbours. We have $d_x\ge d_y$. In the neighbourhoods of
$x,y$ remove each other if they are adjacent. Both sizes drop by the same
amount. $v$ occurs on the $y$ side only. The $x$ side therefore contains
some $z$ missing from the $y$ side, and $z\notin\{v,x,y\}$. Replace edges
$vy,xz$ by $vx,yz$. The four degrees are preserved, simplicity survives,
and the number of desired neighbours increases. In at most $d$ switches the
neighbours of $v$ are precisely $T$. This also proves both rejection tests
necessary. Delete $v$ to obtain the reduced degrees. Conversely unsort using
the stored label permutation and add $v$ adjacent to $T$. Iterate at most
$n$ times. Sorting ties by label makes the construction deterministic.
\end{proof}

\begin{proposition}[Adjacency powers, triangles, and quadrangles]\leavevmode\label{idx:elementary-adjacency-powers-triangles-and-quadrangles}
For a digraph let $A_{uv}$ count labelled arcs from $u$ to $v$.
Then $(A^k)_{uv}$ counts the length-$k$ walks from $u$ to $v$, with
$A^0=I$. For a simple undirected graph,
\[
 \operatorname{tr}(A^2)=2m,\qquad \operatorname{tr}(A^3)=6t,\qquad
 \operatorname{tr}(A^4)=2m+4\sum_v\binom{d(v)}2+8q,
\]
where $t,q$ count unoriented triangles and unoriented simple four-cycles.
\end{proposition}
\begin{proof}
An empty walk is determined by its one vertex. In a walk of length $k+1$
mark the last step $w\to v$. Choosing a prefix to $w$ and one of its last arcs
gives every such walk exactly once. Summing over $w$ gives matrix multiplication.
A closed two-step walk chooses a starting end of an edge. A closed three-step
walk uses three distinct vertices and chooses a start and a direction on its
triangle. A closed four-step walk on two vertices alternates along one edge
(two starts). On three vertices it uses a two-edge path twice, with four
rooted ordered walks for each unordered pair of neighbours of its centre:
start at either leaf, or start at the centre with either first leaf. On four
distinct vertices it chooses one of eight starts with direction. These cases
partition the finite walk table, including walks in graphs containing chords.
For fixed $u$ there are at most $\max(1,m)^k$ labelled directed walks,
and a bounded table of descriptions uses one edge index per step. For a
simple graph one can instead enumerate $k$ subsequent vertices.
\end{proof}

\par\noindent\begin{minipage}{\linewidth}\begin{example}[The six-vertex graph in the source article]\leavevmode\label{ex:six}
Using its labels $1,\ldots,6$, take $12,15,23,25,34,45,46$. The row counts
are $(2,3,2,3,3,1)$ and sum to $14$. Four odd vertices preclude a single
Euler trail. The vertex mask $\{2,4,5\}$ is a cover. The three disjoint edges
$15,23,46$ show that this cover is minimum. The edges
$12,23,34,45,46$ form a spanning tree and the other two edges give two
independent fundamental cycles. There is one triangle ($125$), and hence
$\operatorname{tr}(A^3)=6$. A three-colouring is given by the classes
$\{1,3,6\}$, $\{2,4\}$, $\{5\}$. The triangle certifies optimality.

\begin{center}
\begin{tikzpicture}[x=0.89cm,y=0.89cm,ordinary/.style={circle,draw=ink,fill=white,minimum size=6.6mm,inner sep=0.3pt,font=\small\sffamily},cover/.style={ordinary,fill=ink,text=white}]
\coordinate (a) at (0,0.9); \coordinate (b) at (1.5,1.55); \coordinate (c) at (3.2,1.55);
\coordinate (d) at (3.2,0); \coordinate (e) at (1.5,0); \coordinate (f) at (4.6,-0.05);
\draw[gray] (a)--(b)--(e)--(d)--(c);
\draw[linkcol,line width=1.8pt] (a)--(e) (b)--(c) (d)--(f);
\node[ordinary] at (a){1}; \node[cover] at (b){2}; \node[ordinary] at (c){3};
\node[cover] at (d){4}; \node[cover] at (e){5}; \node[ordinary] at (f){6};
\end{tikzpicture}\par
{\small Filled vertices give the cover. The three heavy edges give the matching.}
\end{center}
\end{example}\end{minipage}\par
